Fix a sufficiently saturated and strongly homogeneous model $\mathbb U$
of $\mathrm{ACVF}(0,0)$. Write $K$ for its valued-field sort,
$\mathcal O$ for the valuation ring, $\mathcal M$ for the maximal ideal,
and $\mathbf k$ for the residue field. We use additive notation for the
value group $\Gamma$. Thus
\[
 \RV^\times=K^\times/(1+\mathcal M),\qquad
 1\longrightarrow\mathbf k^\times\longrightarrow\RV^\times
   \xrightarrow{\vrv}\Gamma\longrightarrow0.
\]
The map $\rv:K\to\RV_\infty$ extends the quotient map by one absorbing
zero, written $\infty$ when it is useful to distinguish it from residue
zero. Definability refers to the original two-sorted RV language: the
ring language on VF, multiplication and totalized inversion on RV,
residue-field addition on its residue domain, the value comparison,
and the map $\rv$. No section of the valuation is named.
For the valued-field and model-theoretic conventions we refer to
\cite{EnglerPrestel,Marker}.

A ground $\mathbb S$ consists of a VF-generated real substructure and
a subgroup of $\Gamma$ containing the values of its RV elements.
We permit parameters in that subgroup as imaginary parameters. The
two parameter sets are smaller than the saturation cardinal, which is
uncountable. We assume throughout that
\[
 \Gamma_0=\dcl^{\mathrm{eq}}(\mathbb S)\cap\Gamma
\]
is nontrivial. Passing to definable closure changes neither the definable
sets nor the definable maps over a ground. It does not mean passing to
algebraic closure. The subgroup $\Gamma_0$ is divisible, and the theory
of nontrivial divisible ordered abelian groups gives
$\Gamma_0\preceq\Gamma$.
For a finite ambient tuple $b$, the notation $\mathbb S_b$ means the
parameter base generated by $\mathbb S$ and $b$, equivalently its
presentation up to definable closure. Every fibre assertion below is
interpreted over this base, with the original Gamma parameters retained.

Unless indicated otherwise, every definable set, including the targets of
the definable maps in the fibre theorems below, consists of real points:
it is a subset of $K^n\times\RV_\infty^r$ for some $n,r$. Gamma-sort
parameters may occur in its definition, but the Gamma sort is imaginary
and contributes no coordinates.

\subsection{Categories, labels and scissors}
An object of $\VF_*$ is a definable subset of
$K^n\times\RV_\infty^r$ whose RV fibre over each VF tuple is finite.
The arities may be zero; no valuative boundedness of its VF coordinates
is imposed. Isomorphisms are definable bijections. An object of
$\RV[k]$ is a pair $\bm U=(U,p)$, where $U$ is a definable RV set and
\[
 p:U\longrightarrow(\RV^\times)^k
\]
is definable and finite-to-one. Isomorphisms need not preserve the
chosen label map, but its degree $k$ is retained before taking the
graded Grothendieck construction. Finite disjoint unions of degrees
are allowed. Zero coordinates and finite tags are retained by finite
stratification. The category $\RES$ consists of RV sets with finite
value support, with definable bijections as isomorphisms; its grading
will be forgotten.

Finiteness of the RV fibres is preserved by ground-definable bijections.
Here is the parameter issue in this assertion. Realize the finitely many
Gamma parameters in a formula for the bijection by RV representatives,
over a small real parameter set. RV exchange over that set implies
invariance of RV-fibre dimension under the bijection
\cite[Lemma 4.13]{YinST}. That dimension is zero precisely when every
RV fibre is finite. Since finiteness concerns the same underlying fibres,
the conclusion holds over the original ground; the auxiliary
representatives do not become parameters of the bijection there.

For any of these definable categories, $K_+$ denotes the positive
Grothendieck semiring, with the relations
$[X]=[Y]+[X\setminus Y]$ for definable subobjects. Multiplication is
Cartesian product. Its group completion is denoted by $K_0$.
Equality of positive classes is equivalent to a definable bijection:
a finite sequence of scissor decompositions can be refined to a common
finite partition, and the piecewise maps assemble on tagged disjoint
unions. Conversely a definable bijection gives equality of classes.
In the group completion one must also allow a common positive
stabilizer; cancellation in the positive semiring is not assumed.

For a parameter base $\mathbb T$ put
\[
 R(\mathbb T)=K_0(\RES_{\mathbb T})/J_{\mathbb T},
\]
where $J_{\mathbb T}$ is generated by
$[\alpha^\sharp]-[\G_m]^n$, for all $\mathbb T$-definable tuples
$\alpha\in\Gamma^n$. Here
\[
 \alpha^\sharp=\{u\in(\RV^\times)^n:\vrv(u)=\alpha\},
 \qquad \G_m=\mathbf k^\times.
\]
The one-coordinate relations generate the same ideal: expand the
difference of two products. Put $g=[\G_m]$ and $L=1+g$ in this ring.
Products with arbitrary residue objects are included in the ideal
relations. All objects, coefficients and bijections defining
$R(\mathbb T)$ are over $\mathbb T$ itself.

\subsection{Lifting and the two retractions}
For $\bm U=(U,p)\in\RV[k]$ define
\[
 \mathbb L\bm U=\{(x,u)\in(K^\times)^k\times U:\rv(x)=p(u)\}.
\]
Over a $(\VF,\Gamma)$-generated ground the lifting theorem identifies
$K_0\VF_*$ with the quotient of $K_0\RV[*]$ by the special relation
$\bm P-1$, where
\[
 \bm P=[\{1\}]-[\rv(\mathcal M\setminus\{0\})]
       \quad\text{in degree }1.
\]
We write $\int$ for the inverse to the induced lifting isomorphism.
At the positive level its existence includes actual lift-bijection
witnesses and the special lifting congruence. We use the lifting and
preparation results of \cite{YinST,YinI} with their original parameter
bases. We do not assume a fibre integral over a base for which the
lifting theorem's generation hypothesis has not been established.

The two Euler characteristics on the value-group sort are denoted by
$\chi_{\mathrm b}$ and $\chi_{\mathrm g}$. The latter is the
o-minimal Euler characteristic of \cite[Chapter 4]{vdD}.
We recall the cylindrical
calculation that will also provide finite first-order certificates.
A cylindrical cell in $\Gamma^{r+1}$ is obtained from a cell in
$\Gamma^r$ by adjoining the graph of a definable function, an open
interval between two ordered definable functions, an upper or lower
open ray, or the whole line. Starting with the one-point cell in
$\Gamma^0$, attach the following factors at each step:
\[
\begin{array}{c|rrrr}
 &\text{graph}&\text{bounded open interval}&\text{open ray}&\text{line}\\ \hline
 \chi_{\mathrm b}&1&-1&0&1\\
 \chi_{\mathrm g}&1&-1&-1&-1.
\end{array}
\]
The value of a cell is the product of its successive factors, and the
value of a finite disjoint cell decomposition is the sum of its cell
values. These formulas agree with successive one-coordinate Euler
integration: a graph contributes one point, a band contributes an open
interval, and the two unbounded cases give the ray and line factors.
Induction on the number of coordinates proves the formula for a
cylindrical cell; finite additivity proves it for a decomposition.
The usual two-coordinate interchange for this integration gives
independence of coordinate order and the Gamma fibre-integration
formula. The bounded convention is equivalently the eventual geometric
Euler characteristic of intersection with a closed box.

For a fixed parameter tuple, divisible ordered-group quantifier
elimination gives a finite cylindrical cell decomposition of every
Gamma-definable fibre. Record each cell as a finite tree of the five
steps just listed, together with formulas for its domain and endpoint
graphs. Record the ordering of the endpoints, the projection identities,
and disjointness and covering of the cells. The Euler value of the tree
is the integer obtained from its shape by the displayed table.
All these validity assertions are first order once the finite trees,
arities and formulas have been fixed. This gives a finite certificate
of the Euler value in the original language, over precisely the fibre's
parameter base. Compactness makes the list of such certificates finite
on a definable parameter family, after a finite definable partition.
In particular the fibre Euler function has finite range and definable
level sets. No common externally chosen Gamma parameters are needed.

\label{sec:retract}
The residue retractions, recalled in \cite[Section 4]{FY}, are characterized
by their monomial evaluations. If a nonzero-RV piece has coordinates
\[
 h_M(U)=V\times I^\sharp,\qquad
 \vrv(V)=\{\beta\},\qquad I\subseteq\Gamma^\ell,
\]
then
\begin{align}
 \mathbb E_{\mathrm b}([\bm U])&=[V]g^\ell\chi_{\mathrm b}(I),
 \label{eq:bounded-evaluation}\\
 \widetilde{\mathbb E}_{\mathrm g}([\bm U])
   &=[V]g^\ell\chi_{\mathrm g}(I).
 \label{eq:geometric-evaluation}
\end{align}
They extend additively across finite zero strata and tags. The second
map is unnormalized; on degree $k$ its normalized version is
$\mathbb E_{\mathrm g}=L^{-k}\widetilde{\mathbb E}_{\mathrm g}$,
with values in $R(\mathbb T)[L^{-1}]$.
The bounded map and the normalized geometric map kill the special
relation, so they apply to $\int[A]$. For example,
$\chi_{\mathrm b}(\Gamma_{>0})=0$ and
$\chi_{\mathrm g}(\Gamma_{>0})=-1$ show that the two values of
$\bm P$ are respectively $1$ and $L^{-1}(1+g)=1$.
Their coefficient and tensor constructions use only definable sets and
maps over the stated base. The monomial formulas express their already
well-defined ring homomorphisms, and do not define a value by an
unproved choice of decomposition.

\subsection{Parameter control and preparation}
We will use two consequences of RV quantifier elimination
\cite{YinQE,YinST}. First, the
residue-field and value-group sorts are orthogonal, with the value group
stably embedded as a pure divisible ordered group. Second, for a finite
nonzero RV tuple $\tau$,
\begin{equation}\label{eq:gamma-control}
 \dcl^{\mathrm{eq}}(\mathbb S,\tau)\cap\Gamma
 =\operatorname{span}_{\Q}
     (\Gamma_0,\vrv(\tau)).
\end{equation}
Here the right side allows rational linear combinations and ground
constants. To see the second assertion, eliminate quantifiers in a formula
uniquely defining a Gamma element. On each leading-term guard the value
conditions are rational affine inequalities whose constants are ground
values and the values of RV monomials in $\tau$. A singleton solution
must be a boundary value of one of these systems; otherwise it lies in
a nontrivial interval in a fibre cell. Solving that boundary equation
gives the asserted rational affine expression. The reverse inclusion
is immediate from divisibility and definability. This also explains why
a finite definable tuple of Gamma parameters may be controlled without
choosing new RV representatives.

The preparation theorem used below has the following form. Given a
finite collection of VF polynomials with ground coefficients, there is
a definable bijection on a VF-definable base, with RV parameters carried
along, whose target is an RV-pullback and on which the RV values of all
the chosen polynomials depend only on the resulting RV tuple.
This is polynomial RV preparation by special bijections
\cite[Theorem 5.5]{YinST}. Apply it after quantifier elimination to the
polynomials occurring in a finite collection of formulas. It then
prepares those formulas simultaneously. Original RV coordinates and
new RV coordinates are retained; an RV parameter is not replaced by
a chosen VF lift. Whole lift-bijection graphs are transported as graphs
of total VF families, rather than asserted to become maps in the RV sort.
